% Move pads through the honest gates, expose the consumed Pauli pair, and twirl.
\newcommand{\mbqcpadgate}[1]{\gate{\textcolor{CVQCTeal}{#1}}}
\newcommand{\mbqcjointattack}{\multigate{1}{\textcolor{CVQCCoral}{A}}}
\newcommand{\mbqcattackghost}{\ghost{\textcolor{CVQCCoral}{A}}}
\begin{frame}{MBQC Analysis}
  \centering
  \begin{tikzpicture}
    \path[use as bounding box] (-6.6,-3.15) rectangle (6.6,3.05);

    \only<1>{
      \node[inner sep=0pt,text=CVQCInk] at (0,0.10) {%
        \resizebox{12.5cm}{!}{%
          \hspace*{2.1cm}%
          $\Qcircuit @C=1em @R=1em {
            \lstick{X^{x_1}T^{t_1}\ket{\psi}} & \ctrl{1} & \gate{HT^\delta}
              & \multigate{1}{\textcolor{CVQCCoral}{A}} & \meter
              & \cgate{\oplus r} & \cw \\
            \lstick{X^{x_2}T^{t_2}\ket{+}} & \control\qw & \qw
              & \ghost{\textcolor{CVQCCoral}{A}} & \qw
              & \qw & \qw
          }$\hspace*{0.4cm}%
        }%
      };
    }
    \only<2>{
      \node[inner sep=0pt,text=CVQCInk] at (0,0.10) {%
        \resizebox{12.5cm}{!}{%
          \hspace*{0.6cm}%
          $\Qcircuit @C=1em @R=1em {
            \lstick{\ket{\psi}} & \gate{\textcolor{CVQCTeal}{X^{x_1}T^{t_1}}}
              & \ctrl{1} & \gate{HT^\delta}
              & \multigate{1}{\textcolor{CVQCCoral}{A}} & \meter
              & \cgate{\oplus r} & \cw \\
            \lstick{\ket{+}} & \gate{\textcolor{CVQCTeal}{X^{x_2}T^{t_2}}}
              & \control\qw & \qw
              & \ghost{\textcolor{CVQCCoral}{A}} & \qw
              & \qw & \qw
          }$\hspace*{0.4cm}%
        }%
      };
    }
    \only<3>{
      \node[inner sep=0pt,text=CVQCInk] at (0,0.10) {%
        \resizebox{12.5cm}{!}{%
          \hspace*{0.6cm}%
          $\Qcircuit @C=1em @R=1em {
            \lstick{\ket{\psi}} & \ctrl{1} & \mbqcpadgate{X^{x_1'}T^{t_1'}}
              & \gate{HT^\delta} & \mbqcjointattack & \meter
              & \cgate{\oplus r} & \cw \\
            \lstick{\ket{+}} & \control\qw & \mbqcpadgate{X^{x_2'}T^{t_2'}}
              & \qw & \mbqcattackghost & \qw & \qw & \qw
          }$\hspace*{0.4cm}%
        }%
      };
    }
    \only<4>{
      \node[inner sep=0pt,text=CVQCInk] at (0,0.10) {%
        \resizebox{12.5cm}{!}{%
          \hspace*{0.6cm}%
          $\Qcircuit @C=1em @R=1em {
            \lstick{\ket{\psi}} & \ctrl{1} & \gate{HT^c}
              & \mbqcpadgate{X^rZ^z} & \mbqcjointattack & \meter
              & \cgate{\oplus r} & \cw \\
            \lstick{\ket{+}} & \control\qw & \mbqcpadgate{X^xT^t}
              & \qw & \mbqcattackghost & \qw & \qw & \qw
          }$\hspace*{0.4cm}%
        }%
      };
    }
    \only<5>{
      \node[inner sep=0pt,text=CVQCInk] at (0,0.10) {%
        \resizebox{12.5cm}{!}{%
          \hspace*{0.6cm}%
          $\Qcircuit @C=1em @R=1em {
            \lstick{\ket{\psi}} & \ctrl{1} & \gate{HT^c}
              & \mbqcpadgate{X^rZ^z} & \mbqcjointattack
              & \mbqcpadgate{(X^rZ^z)^\dagger} & \meter & \cw \\
            \lstick{\ket{+}} & \control\qw & \mbqcpadgate{X^xT^t}
              & \qw & \mbqcattackghost & \qw & \qw & \qw
          }$\hspace*{0.4cm}%
        }%
      };
    }
    \only<6>{
      \node[inner sep=0pt,text=CVQCInk] at (0,0.10) {%
        \resizebox{12.5cm}{!}{%
          \hspace*{0.6cm}%
          $\Qcircuit @C=1em @R=1em {
            \lstick{\ket{\psi}} & \ctrl{1} & \gate{HT^c}
              & \qw & \gate{\textcolor{CVQCCoral}{X^aZ^b}}
              & \qw & \meter & \cw \\
            \lstick{\ket{+}} & \control\qw & \mbqcpadgate{X^xT^t}
              & \qw & \gate{\textcolor{CVQCCoral}{B_Q}}
              & \qw & \qw & \qw
          }$\hspace*{0.4cm}%
        }%
      };
    }
    \only<7>{
      \node[inner sep=0pt,text=CVQCInk] at (0,0.10) {%
        \resizebox{12.5cm}{!}{%
          \hspace*{0.6cm}%
          $\Qcircuit @C=1em @R=1em {
            \lstick{\ket{\psi}} & \ctrl{1} & \gate{HT^c}
              & \qw & \gate{\textcolor{CVQCCoral}{X^aZ^b}}
              & \qw & \meter & \cw \\
            \lstick{\ket{+}} & \control\qw & \mbqcpadgate{X^xT^t}
              & \qw & \gate{\textcolor{CVQCCoral}{B_Q}}
              & \qw & \qw & \rstick{\textcolor{CVQCCoral}{B_Q}X^{x'}T^{t'}\textcolor{CVQCCoral}{X^a}HT^c\ket{\psi}}\qw
          }$\hspace*{3.4cm}%
        }%
      };
    }
  \end{tikzpicture}

  \note<1>{
    \footnotesize
    Return to the encrypted real-copy gadget. The prover may deviate
    arbitrarily. For the analysis, represent its behavior as the prescribed
    unitary followed by a joint attack A immediately before measurement.
    This attack acts on both drawn qubits and can also act on private
    memory, which is omitted from the picture. It may depend on the public
    masked instruction and prior transcript. The manuscript's Honest
    Unitary Before the Attack lemma justifies this placement by absorbing
    earlier deviations into A.
    The verifier still unmasks the report as $\beta=b\oplus r$ and runs
    the key updates from the preceding encrypted-MBQC slide. The surviving
    state may be attacked and correlated with private memory. Keep output
    labels and bookkeeping formulas off this analysis picture.
    Next, commute the pads and apply a partial Pauli twirl to this attack.
    Source: CVQC-paper/4\_Real\_Copy.tex, Honest Unitary Before the Attack
    and Characterization of the State lemmas.
  }
  \note<2>{
    \footnotesize
    Pull each input pad out of the state label and draw it as the first
    gate on its wire. The inputs now read psi and plus, with explicit
    $X^{x_1}T^{t_1}$ and $X^{x_2}T^{t_2}$ pad gates before CZ. This is the
    same circuit as the previous build, with the encryption made visible
    as an operation. The keys are unchanged because the pad gates remain
    before CZ. The coral joint attack, masked measurement, and classical
    XOR correction stay in place. This prepares the next step of commuting
    the pads through CZ and the masked gate for the twirling analysis.
  }
  \note<3>{
    \footnotesize
    Commute the two pad gates to the right of CZ. This is the paper's
    CZ key-update identity, up to global phase. The x bits are unchanged,
    while $t_1'=t_1+4x_2$ and $t_2'=t_2+4x_1$ modulo 8.
    The primes on the pad gates denote these post-CZ keys. The masked
    instruction delta was computed using the post-CZ first-qubit key,
    so its value and the attack A are unchanged. This build exposes the
    plaintext CZ before the updated pads. Explain the algebra verbally.
    Source: Real-Copy Protocol, Update Correctness lemma.
  }
  \note<4>{
    \footnotesize
    Commute the first pad through the masked gate, using the identity
    $HT^\delta X^{x_1'}T^{t_1'}=X^rZ^{x_1'}HT^c$ up to global phase.
    Write z for $x_1'$ and relabel the surviving premeasurement key
    $(x_2',t_2')$ as (x,t). The logical gate is now HT to the c, and the
    measured qubit has a Pauli pad rather than an X-T pad. The lower
    pad remains X to the x times T to the t.
    The Uniformity Preserving lemma supplies independent uniform reference
    coordinates r,z,x,t,delta. In particular, for fixed public delta the
    consumed bits r,z can be averaged without changing the live key
    coordinates. A can depend on delta and the prior transcript; private
    memory remains implicit. Keep carried branch weights throughout.
  }
  \note<5>{
    \footnotesize
    Move the classical XOR-r decoder before measurement as a quantum X-r
    correction. Insert Z-z as well: it changes only a measurement-branch
    phase, leaving the retained density operator unchanged. The result is
    the inverse Pauli pad $(X^rZ^z)^\dagger$ before measurement.
    Thus the joint attack is now between a Pauli pad and its inverse on
    the outgoing qubit. This is the twirling pattern from the previous
    slide. Equality here concerns the surviving ciphertext, prover memory,
    and live verifier keys after the raw report and consumed coordinates
    are discarded. Any raw-outcome copy retained by the prover is included
    in its private memory inside A. Source: Real-Copy Protocol, the
    premeasurement attack decomposition preceding the partial-twirl lemma.
  }
  \note<6>{
    \footnotesize
    Expand the joint attack on the outgoing qubit as
    $A=\sum_Q Q\otimes B_Q$, with $Q=X^aZ^b$. The phase bit b in Q
    is distinct from the raw measurement report, which we call m here.
    Averaging the consumed Pauli key (r,z) removes cross terms between
    different Q labels. The pad and inverse pad disappear from the first
    wire, leaving a Pauli there and its coefficient operator on the
    remaining ciphertext and private memory. The picture is one Kraus
    term in a sum, not a uniformly sampled normalized branch. The lower
    X-T pad stays in place. The possibly nonunitary BQ carries the branch
    weight, as in the CNOT example. No additional scalar probability
    should be assigned to a normalized version of this operator.
  }
  \note<7>{
    \footnotesize
    Reading the teleportation identity with $Q=X^aZ^b$, the X-a component
    flips the decoded outcome; Z-b contributes a branch phase. With decoded
    bit $\beta=m\oplus r$, the lower vector before the final key update
    is proportional to $B_QX^xT^tX^{\beta\oplus a}HT^c\ket{\psi}$.
    Absorb X-beta into the live pad: $x'=x\oplus\beta$ and
    $t'=(-1)^\beta t$ modulo 8. This gives the displayed component
    $B_QX^{x'}T^{t'}X^aHT^c\ket{\psi}$.
    The label uses the paper's positive-state decomposition after summing
    the decoded measurement branches and relabeling the uniform surviving
    keys. Their common factor of one-half is absorbed by that sum; BQ's
    weight remains. For an individual fixed measurement branch, retain
    its one-over-square-root-two amplitude and branch phase. Sum over Q
    and average the reference keys to recover the full joint state.
    Private memory and dependence of BQ on public delta are suppressed.
    Source: Characterization of the State lemma and its proof.
  }
\end{frame}
